% from ega4-14-fr.tex
\oldpage[IV-3]{216}

  ainsi obtenu s'appliquera en particulier à des morphismes \emph{plats} (2.4.6),
  dont l'importance en géométrie algébrique va croissant. On peut considérer que
  les énoncés faisant intervenir l'hypothèse qu'un morphisme est universellement
  ouvert sont des généralisations communes d'énoncés faisant intervenir une
  hypothèse de normalité et d'énoncés faisant intervenir une hypothèse de platitude.
\end{enumerate}
\end{remarks}

\subsection{Morphismes universellement ouverts et quasi-sections}
\label{subsection:IV.14.5-fr}

\begin{lemma}[14.5.1]
\label{IV.14.5.1-fr}
Soient $Y$ un préschéma irréductible localement noethérien, $f:X\to Y$ un
morphisme localement de type fini, $x$ un point de $X$, tel que $f$ soit
équidimensionnel (13.3.2) au point $x$~; on pose $y=f(x)$,
$e=\dim_x(f^{-1}(y))$. Soit $X'$ une partie fermée irréductible de $X$ contenant
$x$, $n$ un entier tel que l'on ait $\operatorname{codim}(X',X)\leq n$ et
$\dim_x(X'\cap f^{-1}(y))\leq e-n$. Alors on a nécessairement
$\operatorname{codim}(X',X)=n$, $\dim_x(X'\cap f^{-1}(y))=e-n$, et la restriction $X'\to Y$
de $f$ est un morphisme équidimensionnel en $x$ (et a fortiori dominant).
\end{lemma}

La question étant locale sur $X$, on peut supposer que $f$ est de type fini et
\emph{équidimensionnel}, de sorte que pour \emph{tout} $z\in Y$, \emph{toutes} les
composantes irréductibles de $f^{-1}(z)$ sont de dimension $e$ (13.3.1, $a''$).
Soient $x'$ le point générique de $X'$, $y'=f(x')$,
$Y'=\overline{\{y'\}}=\overline{f(X')}$ et posons $Z=f^{-1}(Y')$. En vertu de
($\mathbf 0$, 14.2.2), on a
\begin{equation}
  \operatorname{codim}(X',Z)\leq n.
  \tag{14.5.1.1}\label{IV.14.5.1.1-fr}
\end{equation}
Et si les deux membres sont égaux, on a nécessairement $\operatorname{codim}(X',X)=n$ et $Z$
contient une composante irréductible de $X$, ce qui entraîne (13.3.1) que
$f(X')$ est dense dans $Y$ et par suite $Z=X$~; donc l'égalité
$\operatorname{codim}(X',Z)=n$ équivaut à la conjonction de l'égalité $\operatorname{codim}(X',X)=n$ et de
la relation $Z=X$.

D'autre part, en raisonnant dans les préschémas réduits de $Y$ et $X$ ayant $Y'$
et $Z$ respectivement pour espaces sous-jacents, on déduit de (5.1.2) et de
(\textbf{I}, 3.6.5) que l'on a
\begin{equation}
  \operatorname{codim}(X'\cap f^{-1}(y'),f^{-1}(y'))=\operatorname{codim}(X',Z).
  \tag{14.5.1.2}\label{IV.14.5.1.2-fr}
\end{equation}
En vertu de l'hypothèse, $f^{-1}(y')$ est biéquidimensionnel (5.2.1) et de
dimension $e$, donc ($\mathbf 0$, 14.3.5), on a, en vertu de (14.5.1.2) et
(14.5.1.1)
\begin{equation}
  e'=\dim(X'\cap f^{-1}(y'))=e-\operatorname{codim}(X',Z)\geq e-n,
  \tag{14.5.1.3}\label{IV.14.5.1.3-fr}
\end{equation}
l'égalité ayant lieu si et seulement si $\operatorname{codim}(X',X)=n$ et si $f(X')$ est dense
dans $Y$. Enfin, (13.1.1), on a $\dim_x(X'\cap f^{-1}(y))\geq e'$, d'où, par
(14.5.1.3),
\begin{equation}
  \dim_x(X'\cap f^{-1}(y))\geq e'\geq e-n.
  \tag{14.5.1.4}\label{IV.14.5.1.4-fr}
\end{equation}
Or, par hypothèse, on a aussi $\dim_x(X'\cap f^{-1}(y))\leq e-n$, d'où les
conclusions de la proposition.

\begin{corollary}[14.5.2]
\label{IV.14.5.2-fr}
Les hypothèses sur $f$, $X$, $Y$, $x$ étant celles de (14.5.1), supposons en
outre que $x$ ne soit pas point maximal de $f^{-1}(y)$. Alors il existe un
voisinage ouvert affine $U$ de $x$ dans $X$, et une section
$g\in\Gamma(U,\mathcal O_X)$ telle que l'ensemble $X'$ des $x'\in U$ tels que
$g(x')=0$ contienne $x$ et ne contienne aucun point maximal de $f^{-1}(y)$. Pour
tout $g$ ayant ces propriétés, $X'$ est équidimensionnel sur $Y$ au point $x$, et
l'on a
\begin{equation}
  \dim_x(X'\cap f^{-1}(y))=e-1 \qquad\text{et}\qquad \operatorname{codim}(X',X)=1.
  \tag{14.5.2.1}\label{IV.14.5.2.1-fr}
\end{equation}
\end{corollary}

